%% toc: P15 170 170.10
% Simultaneous concentration bands: what the exponential method yields at one sample size, confidence sequences, and their comparison with reused intervals, Bonferroni/Holm allocation, Wald's boundaries and an epoch construction.
\chapter{Simultaneous Concentration Bands and Confidence Sequences}
\label{sec:p15-w02b-p15-063}

An interval computed from the first \(n\) observations is often recomputed
as observations accumulate, and the analyst may stop when the interval
looks decisive. A guarantee that holds at each preselected sample size
does not describe that practice. This chapter states the stronger
property that does, proves that the union-allocation band of
\cref{thm:p15-conc-time-uniform} has it, and compares that band with the
procedures it is most often confused with or replaced by.

\emph{Inputs.} Observations with deterministic known ranges and a common
conditional mean, or more generally bounded martingale differences; a
failure probability \(\delta\in(0,1)\). \emph{Conclusions.} Random
intervals that contain the mean at every sample size simultaneously,
except on an event of probability at most \(\delta\); hence also at every
data-dependent stopping time. The comparisons quantify what is lost or
gained, in validity and in width, relative to a reused fixed-sample
interval, Bonferroni and Holm allocation, Wald's likelihood-ratio
boundaries and an epoch construction built on a maximal inequality.

The chapter begins with the fixed-sample bound itself: what Hoeffding's
exponential method can and cannot give at one sample size, and how his
bounds compare with the Chebyshev and normal benchmarks he used
\cite[pp.~13--15 and~17]{Hoeffding1963BoundedSums}. The fixed-sample
results are those of \cref{sec:p15-bounded-concentration-learning}; the
notation \(S_n\), \(V_n\), \(D(q\Vert p)\), \(g(p)\), \(A_1,A_2,A_3\)
and \(\lambda_*\) is that chapter's.

\section{What the exponential method yields at one sample size}
\label{sec:p15-w02b-p15-063-fixed-sample}

Hoeffding's method bounds an indicator by an exponential and then bounds
the exponential moment. For a sum \(S_n=\sum_{i=1}^n(X_i-EX_i)\), a
threshold \(nt\) and \(h>0\) this is the source's (1.7)--(1.8):
\[
 P(S_n\geq nt)\leq\Psi(h)=e^{-hnt}Ee^{hS_n}.
\]
The function \(\Psi\) depends on the law of the observations. The first
result identifies the smallest quantity that can be extracted from it
uniformly over all laws allowed by the hypotheses.

\begin{hbtheorem}{thm:p15-w02b-p15-063-exponential-optimality}{proved=annex:p15-w02b-p15-063-exponential-optimality}{The divergence bound is the best the exponential method gives}
Assume the hypotheses of \cref{thm:p15-conc-common-range}: independent
\(0\leq X_i\leq1\), \(p=n^{-1}\sum_iEX_i\in(0,1)\) and \(0<t<1-p\).
Define \(\Psi(h)=e^{-hnt}Ee^{hS_n}\) for real \(h\).
\begin{enumerate}
\item For every such law,
 \(\inf_{h\in\mathbb R}\Psi(h)=\inf_{h>0}\Psi(h)\leq e^{-nD(p+t\Vert p)}\).
\item If the \(X_i\) are independent Bernoulli variables of common mean
 \(p\), then
 \(\inf_{h>0}\Psi(h)=e^{-nD(p+t\Vert p)}\), and the infimum is attained
 exactly at \(h=\lambda_*=\ln\{(p+t)(1-p)/[p(1-p-t)]\}\).
\item Consequently, suppose a number \(B=B(n,p,t)\) satisfies
 \(\inf_{h>0}\Psi(h)\leq B\) for every admissible law with average mean
 \(p\). Then \(B\geq e^{-nD(p+t\Vert p)}\). The first bound
 \(A_1\) of \cref{thm:p15-conc-common-range} is the smallest bound,
 depending only on \(n\), \(p\) and \(t\), that optimizing the exponential
 bound over \(h\) can produce.
\item It is not the least upper bound of the probability itself. For
 \(n=1\), the supremum of \(P(X_1-p\geq t)\) over laws on \([0,1]\) with
 mean \(p\) is \(p/(p+t)\); it is attained, and
 \[
             \frac p{p+t}<e^{-D(p+t\Vert p)}.
 \]
\end{enumerate}
Hoeffding states that (2.1) is the best obtainable from (1.7), as the
minimum over \(h\) of its right side under the Bernoulli law (2.5), and
distinguishes this from the best possible bound for the probability,
meaning the least upper bound depending only on the stated quantities
\cite[p.~14 and n.~3; p.~15]{Hoeffding1963BoundedSums}. Part~4 supplies
an explicit case of that distinction.
\end{hbtheorem}

The source compares its bounds with the Chebyshev bounds (1.5)--(1.6)
\cite[p.~13]{Hoeffding1963BoundedSums}. With
\(\sigma^2=n^{-1}\sum_{i=1}^n\operatorname{Var}X_i\), as in its (1.2),
these are the values
\[
 K_5=\frac{\sigma^2}{nt^2}\quad\text{for }P(|\overline X-p|\geq t),
 \qquad
 K_6=\frac{\sigma^2}{\sigma^2+nt^2}\quad\text{for }P(\overline X-p\geq t),
\]
where \(\overline X=n^{-1}\sum_iX_i\). Both follow from
\cref{thm:p15-conc-benchmarks} with \(W=\overline X-p\) and variance
\(\sigma^2/n\). Since \(\overline X-p\geq t\) is the event
\(S_n\geq nt\), \(A_1,A_2,A_3\) bound the same probability as \(K_6\).

\begin{hbtheorem}{thm:p15-w02b-p15-063-chebyshev-comparison}{proved=annex:p15-w02b-p15-063-chebyshev-comparison}{The original comparison with the Chebyshev bounds}
Under the hypotheses of \cref{thm:p15-conc-common-range}, let
\(\sigma^2=n^{-1}\sum_i\operatorname{Var}X_i\).
\begin{enumerate}
\item \(\sigma^2\leq p(1-p)\leq1/4\). Equality \(\sigma^2=p(1-p)\) holds
 exactly when every \(X_i\) takes only the values \(0,1\) almost surely and
 every \(EX_i\) equals \(p\).
\item Suppose \(\sigma^2>0\) and put \(x=nt^2\). There is a unique
 \(x_*=x_*(\sigma^2)>0\) such that \(A_3>K_6\) for \(0<x<x_*\),
 \(A_3=K_6\) at \(x=x_*\), and \(A_3<K_6\) for \(x>x_*\). Moreover
 \[
                       x_*<\frac1{2\sigma^2}-1.
 \]
 For \(nt^2>x_*\), therefore, \(A_1\leq A_2\leq A_3<K_6\leq K_5\).
\item For fixed \(n\), \(p\) and \(\sigma^2>0\) there is \(t_0>0\) such
 that \(K_6<A_1\) for \(0<t<t_0\). Even the first bound is larger than
 the one-sided Chebyshev bound for sufficiently small deviations.
\item For fixed \(t\), \(A_1\leq A_2\leq A_3=e^{-2nt^2}\) decreases
 exponentially in \(n\). If the variances satisfy \(\sigma^2\geq s>0\)
 along a sequence of sample sizes, then \(K_6\geq s/(s+nt^2)\), which is
 not smaller than order \(1/n\), and \(A_3/K_6\to0\).
\item For the two-sided event, \(P(|\overline X-p|\geq t)\leq2A_3\) for
 every \(t>0\), by \cref{thm:p15-conc-unequal-ranges} with \(a_i=0\),
 \(b_i=1\) at the threshold \(nt\). Suppose \(\sigma^2>0\) and put
 \(x=nt^2\). Then \(2A_3<K_5\) exactly when \(2xe^{-2x}<\sigma^2\). The
 set of \(x>0\) with \(2A_3\geq K_5\) is a closed interval
 \([x_-,x_+]\), where
 \[
            0<x_-<\tfrac12\ln2<\tfrac12<x_+<\ln(2/\sigma^2).
 \]
 For \(x<x_-\) both \(2A_3\) and \(K_5\) exceed one, so neither bound
 says anything; for \(x>x_+\), \(2A_3<K_5\).
\end{enumerate}
The source summarizes this as an improvement over (1.5)--(1.6) except for
small \(t\) or small \(n\), with exponential decay in \(n\) for fixed
\(t\) \cite[p.~14]{Hoeffding1963BoundedSums}. Parts~2--5 state exactly
where that summary holds for the displayed bounds: the one-sided and
two-sided thresholds depend on the product \(nt^2\) and on the variance.
\end{hbtheorem}

The comparison of Hoeffding's three bounds with one another, stated and
proved as \cref{thm:p15-conc-common-comparison}, is the source's
Section~3 \cite[pp.~18--19]{Hoeffding1963BoundedSums}. There the bounds
are written as \(A_1\leq A_2\leq A_3\) in (3.1)--(3.4), with the remark
that \(A_2=A_3\) exactly when \(\mu=1/2\); the two-sided estimates of
\(A_1/A_2\) for \(\mu>1/2\) and \(\mu=1/2\) are (3.6) and (3.7); for
\(\mu<1/2\) the identity (3.8) gives \(A_1=A_2\) at \(t=1-2\mu\), and
(3.9) is the lower estimate \(A_1/A_2\geq e^{-n(1-2\mu-t)^2/\mu}\). The
same page observes that, when \(t/\mu\) is small enough for the leading
term of the expansion (3.5) to dominate, \(A_1/A_2\) is near one when
\([1/\{2\mu(1-\mu)\}-g(\mu)]t^2n\) is small; the limit statement in
\cref{thm:p15-conc-common-comparison} makes that remark exact. The
source's \(\mu\) is the average mean written \(p\) in
\cref{sec:p15-bounded-concentration-learning}.

For a normally distributed mean the same exponential method gives the
source's (2.16):
if \(\overline X\) is normal with mean \(\mu\) and variance \(\sigma^2/n\),
then
\(P(\overline X-\mu\geq t)=\Phi(-\sqrt nt/\sigma)\leq e^{-nt^2/(2\sigma^2)}\)
\cite[p.~17]{Hoeffding1963BoundedSums}.

\begin{hbtheorem}{thm:p15-w02b-p15-063-normal-benchmark}{proved=annex:p15-w02b-p15-063-normal-benchmark}{Coincidences with the normal-model bound and where it is crude}
Write \(\Phi\) and \(\phi\) for the standard normal distribution function
and density.
\begin{enumerate}
\item If \(S_n\) is normal with mean zero and variance \(n\sigma^2>0\), then
 for \(t>0\)
 \[
       \inf_{h>0}e^{-hnt}Ee^{hS_n}=e^{-nt^2/(2\sigma^2)}.
 \]
 Hence \(\Phi(-y)\leq e^{-y^2/2}\) for \(y>0\).
\item Under the hypotheses of \cref{thm:p15-conc-common-range} with
 \(p\geq1/2\) and \(\sigma^2>0\),
 \[
       A_2=e^{-nt^2/(2p(1-p))}\geq e^{-nt^2/(2\sigma^2)},
 \]
 with equality exactly in the binomial case: every \(X_i\in\{0,1\}\)
 almost surely with \(EX_i=p\).
\item For independent \(a_i\leq X_i\leq b_i\) with \(V_n>0\) and
 \(\sigma^2>0\), the bound of \cref{thm:p15-conc-unequal-ranges} at the
 threshold \(nt\) satisfies
 \[
       e^{-2n^2t^2/V_n}\geq e^{-nt^2/(2\sigma^2)},
 \]
 with equality exactly when every \(X_i\) with \(a_i<b_i\) takes the values
 \(a_i\) and \(b_i\) with probability \(1/2\) each.
\item For every \(y>0\), \(e^{-y^2/2}>\Phi(-y)\) and
 \[
       \frac{e^{-y^2/2}}{\Phi(-y)}\geq\sqrt{2\pi}\,y .
 \]
 The normal-model bound is never attained and its relative error is
 unbounded as \(y\to\infty\). By contrast, \(A_1\) is attained at the
 largest nontrivial deviation \(t=1-p\) (\cref{thm:p15-conc-common-range}),
 and the variance bound \(B_1\) at \(t=b\)
 (\cref{thm:p15-variance-optimality-boundary}).
\item Under \cref{def:p15-upper-variance-setting} with \(\sigma^2>0\) and
 \(0<t<b\), let \(B_1,B_2,\tau,\lambda,h_2\) be as in
 \cref{thm:p15-variance-tail-chain} and \(N=e^{-nt^2/(2\sigma^2)}\). Then
 \[
  e^{-nt^3/(2b\sigma^2)}\leq\frac{B_1}N\leq\frac{B_2}N
   \leq\frac{e^{-\tau h_2(\lambda)}}N\leq e^{nbt^3/(6\sigma^4)} .
 \]
 So Bennett's and Bernstein's bounds and the variance bound \(B_1\) all
 differ from the normal-model value by factors tending to one when
 \(nt^3\to0\) with \(b,\sigma^2\) fixed; in particular for fixed \(n\) as
 \(\lambda=bt/\sigma^2\to0\).
\end{enumerate}
Parts~2--4 are the source's remarks that (2.2) with \(\mu\geq1/2\)
equals the right side of (2.16) in the binomial case (2.5), that (2.6)
equals it for symmetric two-point laws, that (2.1) and (2.8) are attained
at the largest nontrivial \(t\), and that the normal-based bound is crude
for large \(y\) \cite[p.~17]{Hoeffding1963BoundedSums}. Part~5 makes
precise the same page's remark that Bernstein's (2.13), and likewise the
bounds of Theorem~3, are close to the right side of (2.16) when
\(\lambda\) is small. The inequalities in parts~2 and~3 add that the bounded-variable bounds are never smaller
than the normal-model value with the actual variance.
\end{hbtheorem}

\begin{hbexample}{def:p15-w02b-p15-063-chebyshev-example}{Where the Chebyshev bound is smaller}
Take independent fair binary variables, so \(p=1/2\) and
\(\sigma^2=1/4\), the equality case of part~1 of
\cref{thm:p15-w02b-p15-063-chebyshev-comparison}. The function
\(\tfrac14e^{2x}-\tfrac14-x\) has its positive zero at
\(x_*\approx0.6282\), below the guaranteed upper limit
\(1/(2\sigma^2)-1=1\). At \(t=0.1\) this means \(A_3<K_6\) exactly when
\(n\geq63\). Numerically, at \(n=62\) one has
\(A_1\approx0.28696\), \(A_3\approx0.28938\) and \(K_6\approx0.28736\):
the first bound is already below the Chebyshev value, the third not yet.
At \(n=200\), \(A_3\approx0.0183\) while \(K_6=1/9\approx0.111\). At
\(n=10\), \(K_6\approx0.714\) is smaller than \(A_1\approx0.818\).
For the two-sided event, part~5 gives \(x_-\approx0.1787\) and
\(x_+\approx1.0766\). At \(t=0.1\), the bound \(2A_3\) is smaller than
\(K_5\) for \(n\leq17\), where both exceed one, and again exactly from
\(n=108\) on: at \(n=107\), \(2A_3\approx0.2353\) and
\(K_5\approx0.2336\); at \(n=108\), \(2A_3\approx0.2307\) and
\(K_5\approx0.2315\). The two-sided comparison thus needs a larger sample
than the one-sided one before the exponential bound wins.
\end{hbexample}

\section{Confidence sequences}
\label{sec:p15-w02b-p15-063-confidence-sequences}

\begin{hbdefinition}{def:p15-w02b-p15-063-confidence-sequence}{Confidence sequences and random times}
Let \((\Omega,\mathcal F)\) carry a filtration
\((\mathcal F_n)_{n\geq0}\), let \(\mathcal P\) be a family of probability
laws on \(\mathcal F\), and let \(\theta:\mathcal P\to\mathbb R\) be the
target parameter. For each \(n\geq1\) let \(L_n\leq U_n\) be
\(\mathcal F_n\)-measurable extended-real random variables and
\(\mathcal C_n=[L_n,U_n]\cap\mathbb R\).
For \(0<\delta<1\), \((\mathcal C_n)_{n\geq1}\) is a
\emph{\((1-\delta)\) confidence sequence} for \(\theta\) over
\(\mathcal P\) if, for every \(P\in\mathcal P\),
\[
 P\{\theta(P)\in\mathcal C_n\ \text{for every }n\geq1\}\geq1-\delta .
\]
The event is a countable intersection of the events
\(\{L_n\leq\theta(P)\leq U_n\}\in\mathcal F_n\).

A \emph{random time} is a measurable map
\(\tau:\Omega\to\{1,2,\ldots\}\cup\{\infty\}\); it is a \emph{stopping
time} if \(\{\tau=n\}\in\mathcal F_n\) for every \(n\). On
\(\{\tau=n\}\) put \(\mathcal C_\tau=\mathcal C_n\). Coverage at
\(\tau\) is the probability of \(\{\tau<\infty,\ \theta(P)\in\mathcal
C_\tau\}\). \emph{Fixed-time coverage} at level \(1-\delta\) means only
\(P\{\theta(P)\in\mathcal C_n\}\geq1-\delta\) for each separate \(n\).
\index{confidence sequence}\index{coverage!simultaneous}
\end{hbdefinition}

\begin{hbtheorem}{thm:p15-w02b-p15-063-cs-equivalence}{proved=annex:p15-w02b-p15-063-cs-equivalence}{Simultaneous coverage is coverage at every stopping time}
Fix \(P\in\mathcal P\), write \(\theta=\theta(P)\), and let
\(c=P\{\theta\in\mathcal C_n\text{ for every }n\geq1\}\). Then
\[
 c=\inf_\tau P\{\tau<\infty,\ \theta\in\mathcal C_\tau\},
\]
where the infimum may range over any one of these classes: all random
times with \(P(\tau<\infty)=1\); all such stopping times; all bounded
stopping times. In particular the following are equivalent:
simultaneous coverage at level \(1-\delta\); coverage at least \(1-\delta\)
at every almost surely finite random time; the same at every almost surely
finite stopping time; the same at every bounded stopping time.
The running intersections \(\mathcal C'_n=\bigcap_{m\leq n}\mathcal C_m\)
have the same simultaneous coverage event, so they form a
\((1-\delta)\) confidence sequence whenever \((\mathcal C_n)\) does; they
may be empty, but only on the failure event.
Fixed-time coverage at every \(n\) does not imply simultaneous coverage
(\cref{thm:p15-w02b-p15-063-reuse-failure}).
\end{hbtheorem}

\begin{hbtheorem}{thm:p15-w02b-p15-063-union-band-cs}{proved=annex:p15-w02b-p15-063-union-band-cs}{The union-allocation band is a confidence sequence}
Fix deterministic finite numbers \(a_i\leq b_i\), \(i\geq1\), and put
\(V_n=\sum_{i=1}^n(b_i-a_i)^2\). Let \(\mathcal P\) consist of the laws
under which \(X_i\) is \(\mathcal F_i\)-measurable,
\(a_i\leq X_i\leq b_i\) almost surely, and
\(E(X_i\mid\mathcal F_{i-1})=\mu(P)\) for every \(i\geq1\), with a
constant \(\mu(P)\). For \(0<\delta<1\) put
\[
 \rho_n=\sqrt{\frac{V_n}{2}\ln\frac{2n(n+1)}\delta},\qquad
 \mathcal C_n=\left[\overline X_n-\frac{\rho_n}n,\
                     \overline X_n+\frac{\rho_n}n\right],
\]
with \(\overline X_n=n^{-1}\sum_{i\leq n}X_i\) and \(\rho_n=0\) when
\(V_n=0\).
\begin{enumerate}
\item \((\mathcal C_n)\) is a \((1-\delta)\) confidence sequence for
 \(\mu\) over \(\mathcal P\). Independent identically distributed
 observations with values in \([a,b]\) belong to \(\mathcal P\) with
 their natural filtration.
\item For \(a_i=0,b_i=1\), \(\rho_n/n=r_n\), where
 \[
       r_n=\sqrt{\frac{\ln(2n(n+1)/\delta)}{2n}},
 \]
 and the clipped intervals
 \([\max\{0,\overline X_n-r_n\},\min\{1,\overline X_n+r_n\}]\) are again a
 \((1-\delta)\) confidence sequence for \(\mu\).
\item If instead \(E(X_i\mid\mathcal F_{i-1})=\mu_i\) with deterministic
 but unequal \(\mu_i\), the same intervals cover the running average
 \(\overline\mu_n=n^{-1}\sum_{i\leq n}\mu_i\) for every \(n\)
 simultaneously, with probability at least \(1-\delta\). They do not then
 cover a single fixed parameter.
\end{enumerate}
\end{hbtheorem}

\section{Repeating a fixed-sample interval}
\label{sec:p15-w02b-p15-063-reuse}

\begin{hbdefinition}{def:p15-w02b-p15-063-reused-intervals}{The reused interval and the finite-horizon allocation}
In the setting of \cref{thm:p15-w02b-p15-063-union-band-cs} with
\(a_i=0,b_i=1\), put
\[
 f_n=\sqrt{\frac{\ln(2/\delta)}{2n}},\qquad
 \mathcal I^{\rm fix}_n=[\overline X_n-f_n,\ \overline X_n+f_n].
\]
This is the two-sided fixed-sample Hoeffding interval for sample size
\(n\), recomputed at every \(n\). For a horizon \(N\geq1\) and
\(1\leq n\leq N\), the finite-horizon allocation uses
\[
 u_n^{(N)}=\sqrt{\frac{\ln(2N/\delta)}{2n}},\qquad
 \mathcal B^{(N)}_n=[\overline X_n-u^{(N)}_n,\ \overline X_n+u^{(N)}_n].
\]
\end{hbdefinition}

\begin{hbtheorem}{thm:p15-w02b-p15-063-reuse-failure}{proved=annex:p15-w02b-p15-063-reuse-failure}{Reusing a fixed-sample interval fails on overlapping means}
In the model of \cref{def:p15-w02b-p15-063-reused-intervals} and for
\(0<\delta<1\):
\begin{enumerate}
\item \(P\{\mu\in\mathcal I^{\rm fix}_n\}\geq1-\delta\) for each fixed
 \(n\).
\item If the \(X_i\) are independent and identically distributed with
 variance \(\sigma^2>0\), then almost surely
 \(\mu\notin\mathcal I^{\rm fix}_n\) for infinitely many \(n\). Thus
 \(P\{\mu\in\mathcal I^{\rm fix}_n\text{ for every }n\}=0\), and for every
 \(\varepsilon>0\) there is a bounded stopping time \(\tau\) with
 \(P\{\mu\in\mathcal I^{\rm fix}_\tau\}<\varepsilon\). One such time stops
 at the first \(n\leq N\) whose interval excludes \(\mu\), and otherwise
 at \(N\), for \(N\) large.
\item If \(\sigma^2=0\), every \(X_i=\mu\) almost surely and every
 interval covers \(\mu\).
\item \(P\{\mu\in\mathcal B^{(N)}_n\text{ for }1\leq n\leq N\}\geq1-\delta\);
 nothing is asserted for \(n>N\).
\item For independent Bernoulli variables of mean \(p\), let
 \[
       \mathsf A_m=\{k\in\{0,\ldots,m\}:|k/m-p|\leq f_m\},
 \]
 \(v_0(0)=1\), and for \(m\geq1\)
 \[
       v_m(k)=\mathbf1_{\mathsf A_m}(k)\,
              \bigl[p\,v_{m-1}(k-1)+(1-p)\,v_{m-1}(k)\bigr],
 \]
 with \(v_{m-1}(j)=0\) for \(j\notin\{0,\ldots,m-1\}\). Then
 \(v_N(k)=P\{\sum_{i\leq N}X_i=k,\ p\in\mathcal I^{\rm fix}_m
 \text{ for }m\leq N\}\), and the probability of some exclusion by time
 \(N\) is \(1-\sum_kv_N(k)\).
\end{enumerate}
The successive means share their observations, so the exclusion events
are dependent. The formula \(1-(1-\alpha)^M\) of
\cref{thm:p15-se-repeated-fixed} is for independent inspections and does
not apply; part~2 settles the unbounded case without it, and part~5
gives the exact finite-horizon value.
\end{hbtheorem}

\begin{hbexample}{def:p15-w02b-p15-063-reuse-example}{How fast the error grows for fair coin means}
Take fair binary observations and \(\delta=0.05\). Evaluating the
recursion of part~5 gives the probability that some reused interval
\(\mathcal I^{\rm fix}_m\), \(m\leq N\), excludes \(1/2\): about
\(0.0530\) for \(N=100\) and about \(0.1032\) for \(N=1000\). Both
exceed \(\delta\), although every individual interval has coverage at
least \(0.95\). The growth is much slower than for independent
inspections, where \(1-0.95^{100}\approx0.9941\), because overlapping
means are strongly positively dependent; part~2 shows that it still
reaches one. The same recursion with \(r_m\) in place of \(f_m\) gives an
exclusion probability of about \(1.0\times10^{-5}\) by \(N=100\) and
\(1.2\times10^{-5}\) by \(N=1000\), well below the guaranteed \(0.05\).
\end{hbexample}

\section{Error allocation across times: Bonferroni and Holm}
\label{sec:p15-w02b-p15-063-allocation}

For a candidate value \(\mu_0\in[0,1]\) and \(n\geq1\), put
\[
 \pi_n(\mu_0)=\min\{1,\ 2\exp(-2n(\overline X_n-\mu_0)^2)\},
 \qquad w_n=\frac1{n(n+1)} .
\]
Each \(\pi_n(\mu_0)\) is a time-indexed obtained level for the hypothesis
\(\mu=\mu_0\), in the sense of \cref{def:p15-mt-obtained-levels}.

\begin{hbtheorem}{thm:p15-w02b-p15-063-allocation}{proved=annex:p15-w02b-p15-063-allocation}{The time index as a family of hypotheses}
In the model of \cref{thm:p15-w02b-p15-063-union-band-cs} with
\(a_i=0,b_i=1\), and for \(0<\delta<1\):
\begin{enumerate}
\item If \(\mu(P)=\mu_0\), then \(P\{\pi_n(\mu_0)\leq u\}\leq u\) for every
 \(n\geq1\) and \(0\leq u\leq1\).
\item For \(\mu_0\in[0,1]\): \(\mu_0\notin\mathcal C_n\) exactly when
 \(\pi_n(\mu_0)<\delta w_n\); \(\mu_0\notin\mathcal I^{\rm fix}_n\)
 exactly when \(\pi_n(\mu_0)<\delta\); and, for \(n\leq N\),
 \(\mu_0\notin\mathcal B^{(N)}_n\) exactly when
 \(\pi_n(\mu_0)<\delta/N\).
\item Hence the union band is a weighted Bonferroni allocation over the
 countable family of time-indexed tests, with weights summing to one; the
 reused interval is unadjusted testing at every time; the finite-horizon
 band is equal-weight Bonferroni over \(N\) times. The coverage of the
 first and third uses only part~1 and subadditivity, with no assumption on
 the joint dependence of the \(\pi_n\).
\item Fix \(N\geq1\) and \(\mu_0\), and regard \(\pi_1(\mu_0),\ldots,
 \pi_N(\mu_0)\) as obtained levels for \(N\) hypotheses, each of which is
 the statement \(\mu=\mu_0\). Holm's procedure at level \(\delta\)
 (\cref{def:p15-mt-holm-rule}) rejects at least one of them exactly when
 \(\min_{n\leq N}\pi_n(\mu_0)\leq\delta/N\), which is exactly when
 Bonferroni at level \(\delta\) rejects at least one. With fixed positive
 weights \(c_n\), the weighted procedure of \cref{def:p15-mt-weighted}
 rejects at least one exactly when some
 \(\pi_n(\mu_0)\leq\delta c_n/\sum_{\ell\leq N}c_\ell\). The set of values
 \(\mu_0\) rejected by no member is therefore the same for the step-down
 and the single-step procedure.
\item \(u^{(N)}_n<r_n\) exactly when \(N<n(n+1)\), and
 \(u^{(N)}_n=r_n\) exactly when \(N=n(n+1)\).
\end{enumerate}
\end{hbtheorem}

Holm's gain over Bonferroni comes from rejecting distinct hypotheses
one after another (\cref{thm:p15-mt-domination}). Here every member of the
family is the same statement about \(\mu\); once one member is rejected,
the remaining comparisons can only reject further copies of it. Holm's
ordering also uses all \(N\) obtained levels together, so it supplies no
decision at an intermediate time \(n<N\). The countable weighting is the
same device as the fixed hypothesis weights in
\cref{thm:p15-conc-weighted-learning}, applied to times instead of
candidate predictors.

\section{Wald's boundaries as a time-uniform test}
\label{sec:p15-w02b-p15-063-wald}

\begin{hbtheorem}{thm:p15-w02b-p15-063-sprt-comparison}{proved=annex:p15-w02b-p15-063-sprt-comparison}{Wald's likelihood boundaries compared with the union band}
Let \(0<p_0<p_1<1\) and let the observations be independent Bernoulli
variables of mean \(p\) under \(P_p\). Use the likelihood ratio \(L_m\)
and log ratio \(Z_m\) of \cref{def:p15-se-model} for
\(P_0=\operatorname{Bernoulli}(p_0)\) and
\(P_1=\operatorname{Bernoulli}(p_1)\), the score values \(z_+,z_-\) of
\cref{thm:p15-se-binomial}, and put \(\Delta=z_+-z_->0\). Fix \(A>1\),
\(a=\ln A\).
\begin{enumerate}
\item \(P_{p_0}\{Z_m\geq a\text{ for some }m\geq1\}\leq1/A\). With
 \(A=1/\delta\), the upper boundary alone tests \(p_0\) with error at most
 \(\delta\) over unlimited inspection, without any allocation over times.
\item For every \(m\geq1\),
 \(Z_m\geq a\) exactly when
 \(\overline X_m\geq c_*+a/(m\Delta)\), where \(c_*=-z_-/\Delta\) and
 \(p_0<c_*<p_1\).
\item If \(p_0<p<c_*\), there is a unique \(h>0\) with
 \(pe^{hz_+}+(1-p)e^{hz_-}=1\), and
 \(P_p\{Z_m\geq a\text{ for some }m\}\leq A^{-h}<1\). Under the same
 \(P_p\), for every \(\delta\in(0,1)\), almost surely the union band
 \(\mathcal C_n\) excludes \(p_0\) for all sufficiently large \(n\).
\item If \(p>c_*\), then \(P_p\{Z_m\geq a\text{ for some }m\}=1\).
\end{enumerate}
For the two-boundary rule of \cref{def:p15-se-sprt}, part~1 is the
one-sided counterpart of \(\alpha'\leq1/A\) in
\cref{thm:p15-se-errors}.
\end{hbtheorem}

Both constructions control an error over every sample size. Wald's
boundary does it with one nonnegative martingale and pays no logarithmic
price in \(m\), but it answers a different question: on the mean scale its
boundary approaches the fixed point \(c_*\) between the two simple
hypotheses, not the true mean. Part~3 shows that it may never reject
\(p_0\) when the truth lies between \(p_0\) and \(c_*\), while the
union band rejects every false value eventually. The two-boundary test
also stops and accepts one of two simple laws; a confidence sequence makes
no terminal decision and is defined for every \(n\).

\section{An epoch construction from the maximal inequality}
\label{sec:p15-w02b-p15-063-epochs}

The union band applies a fixed-time bound at every \(n\). The maximal
inequality of \cref{thm:p15-conc-martingale}, which rests on
\cref{thm:p15-doob-weak-maximal}, controls a whole range of times at the
price of one. Grouping times into epochs of doubling length gives
another band. For \(k\geq0\) put \(N_k=2^{k+1}-1\); epoch \(k\) is
\(\{2^k,\ldots,N_k\}\), and \(k(n)=\lfloor\log_2n\rfloor\) is the epoch
of \(n\).

\begin{hbtheorem}{thm:p15-w02b-p15-063-epoch-band}{proved=annex:p15-w02b-p15-063-epoch-band}{A confidence sequence from maximal inequalities over epochs}
In the setting of \cref{thm:p15-conc-martingale}, with deterministic
bounds for every \(i\geq1\), and for \(0<\delta<1\), put \(k=k(n)\) and
\[
 \varepsilon_n=\sqrt{\frac{V_{N_k}}2\ln\frac{2(k+1)(k+2)}\delta},
\]
replaced by \(\varepsilon_n=0\) when \(V_n=0\). Then
\[
 P\{|M_n|\leq\varepsilon_n\ \text{for every }n\geq1\}\geq1-\delta.
\]
In the model of \cref{thm:p15-w02b-p15-063-union-band-cs} with
\(a_i=0,b_i=1\), the intervals
\([\overline X_n-e_n,\overline X_n+e_n]\), clipped to \([0,1]\), with
\[
 e_n=\frac1n\sqrt{\frac{N_k}2\ln\frac{2(k+1)(k+2)}\delta},
\]
form a \((1-\delta)\) confidence sequence for \(\mu\).
\end{hbtheorem}

\begin{hbtheorem}{thm:p15-w02b-p15-063-width-comparison}{proved=annex:p15-w02b-p15-063-width-comparison}{Widths of the union, epoch and fixed-sample radii}
For \(a_i=0,b_i=1\), \(0<\delta<1\) and \(k=k(n)\),
\[
 \frac{e_n^2}{r_n^2}
  =\frac{N_k}n\cdot\frac{\ln(2(k+1)(k+2)/\delta)}{\ln(2n(n+1)/\delta)},
 \qquad
 \frac{r_n^2}{f_n^2}=\frac{\ln(2n(n+1)/\delta)}{\ln(2/\delta)}>1 .
\]
\begin{enumerate}
\item At \(n=1\), \(e_1=r_1=\sqrt{\ln(4/\delta)/2}\) for every \(\delta\).
 The clipped interval at \(n=1\) is all of \([0,1]\) for every value
 of \(X_1\) exactly when
 \(\delta\leq4e^{-2}\approx0.5413\).
\item At \(n=2\), \(e_2=\sqrt{3/2}\,r_2\) for every \(\delta\). At every
 epoch start \(n=2^k\), \(k\geq1\), \(e_n>r_n\) for all sufficiently
 small \(\delta\). At every epoch end \(n=N_k\), \(k\geq1\), \(e_n<r_n\) for every \(\delta\).
\item Let \(k_0(\delta)\) be the least \(k\geq2\) with
 \(\delta\,2^k(2^k+1)\geq2(k+1)^2(k+2)^2\). Then \(e_n<r_n\) for every
 \(n\geq2^{k_0(\delta)}\), and \(e_n/r_n\to0\) as \(n\to\infty\).
 One has \(k_0(0.05)=10\), and \(k_0(\delta)=7\) for
 \(10368/16512\leq\delta<1\).
\item As \(\delta\uparrow1\), the radii decrease to the positive limits
 \(\sqrt{\ln(2n(n+1))/(2n)}\),
 \(n^{-1}\sqrt{N_k\ln(2(k+1)(k+2))/2}\) and \(\sqrt{\ln2/(2n)}\).
 No band shrinks to a point as the failure probability approaches one.
\item In the general martingale form, \(\rho_n=0\) exactly when
 \(V_n=0\). Before the zero replacement, the epoch radius at such an
 \(n\) is positive whenever \(V_{N_k}>0\); the replacement changes the
 band only on the null event \(\{V_n=0,\ M_n\ne0\}\). With unequal widths
 \(\varepsilon_n^2/\rho_n^2\) carries the factor \(V_{N_k}/V_n\), which
 can be unbounded: for widths \(b_i-a_i=2^i\), the ratio
 \(\varepsilon_n/\rho_n\) at \(n=2^k\) tends to infinity with \(k\).
\end{enumerate}
\end{hbtheorem}

\begin{hbexample}{def:p15-w02b-p15-063-width-table}{The radii at the conventional level}
At \(\delta=0.05\), for observations in \([0,1]\), the four radii are
(rounded to four decimals):
\[
\begin{array}{r|cccc}
 n & r_n & e_n & f_n & u^{(1000)}_n\\ \hline
 1 & 1.4802 & 1.4802 & 1.3581 & 2.3018\\
 2 & 1.1705 & 1.4336 & 0.9603 & 1.6276\\
 3 & 1.0144 & 0.9557 & 0.7841 & 1.3289\\
 4 & 0.9141 & 1.1621 & 0.6791 & 1.1509\\
 10 & 0.6477 & 0.7081 & 0.4295 & 0.7279\\
 100 & 0.2541 & 0.2213 & 0.1358 & 0.2302\\
 1000 & 0.0936 & 0.0655 & 0.0429 & 0.0728\\
 10^4 & 0.0332 & 0.0272 & 0.0136 & \text{none}\\
 10^6 & 0.0040 & 0.0023 & 0.0014 & \text{none}
\end{array}
\]
The column \(f_n\) is not a confidence sequence
(\cref{thm:p15-w02b-p15-063-reuse-failure}); it shows the price of
simultaneity. The finite-horizon column covers only \(n\leq1000\); it is
wider than \(r_n\) while \(n(n+1)<1000\) and narrower afterwards. The
epoch radius alternates within early epochs. Direct evaluation of the two
closed forms shows that \(n=528\) is the last sample size with
\(e_n\geq r_n\); the guarantee of part~3 begins at \(2^{10}=1024\). With
\(n=1\) every band is the whole unit interval at this level.

For Wald's boundaries take \(p_0=1/2\), \(p_1=3/4\) and \(A=20\). Then
\(c_*=\ln2/\ln3\approx0.6309\). The upper boundary on the mean scale is
about \(0.9036\), \(0.6582\) and \(0.6337\) at \(m=10,100,1000\); the
union band at \(\delta=0.05\) excludes \(1/2\) from above once the mean
exceeds about \(0.7541\), \(0.5936\) and \(0.5332\) at
\(m=100,1000,10^4\). For true mean \(p=0.55\), the root of part~3 is
\(h\approx0.6200\), so Wald's upper boundary is ever crossed with
probability at most \(20^{-h}\approx0.156\), while the union band excludes
\(1/2\) eventually with probability one.
\end{hbexample}

\begin{hbwarning}{def:p15-w02b-p15-063-limitations}{What the simultaneous bands do not provide}
The bands require ranges fixed before sampling and a conditional mean that
does not change; with changing deterministic conditional means they cover
the running average instead
(\cref{thm:p15-w02b-p15-063-union-band-cs}). A band for one prespecified
parameter does not cover a parameter selected after looking at several;
that selection needs its own allocation, as in
\cref{thm:p15-conc-finite-learning}. The coverage probability is a
sampling probability under every law of the model, not a posterior
probability for \(\mu\). At conventional levels the first few intervals
are uninformative (part~1 of \cref{thm:p15-w02b-p15-063-width-comparison}).
None of the comparisons above shows that either band has the smallest
possible width among confidence sequences; they compare the displayed
constructions only.
\end{hbwarning}

\section{Choosing and reporting a simultaneous band}
\label{sec:p15-w02b-p15-063-method}

\begin{hbmethod}{def:p15-w02b-p15-063-method}{From an inspection plan to a valid band}
\begin{enumerate}
\item Decide whether the interval will be read at one preselected sample
 size, at a fixed finite set of sizes, or at sizes chosen while observing.
 Only the first case permits the fixed-sample radius \(f_n\).
\item Verify the model: deterministic ranges \(a_i\leq X_i\leq b_i\) known
 before sampling, and a common conditional mean given the past.
 Independence is not required.
\item For a fixed finite set of \(N\) inspection times, the equal-weight
 allocation \(u^{(N)}_n\) is valid at those times only. For inspection at
 data-dependent times without a fixed horizon, use \(r_n\) or \(e_n\), or
 any other countable allocation with weights summing to at most one.
\item Choose the allocation before observing: \(r_n\) is never wider than
 \(e_n\) at \(n=1\) and is narrower at epoch starts for small \(\delta\);
 \(e_n\) is narrower for large \(n\). Taking the smaller of the two at
 each \(n\) after seeing the data is not covered by either guarantee; an
 intersection of both bands, each at level \(\delta/2\), is.
\item Report the level, the allocation, the range assumption and the
 stopping rule. Clip the interval to the known range of the mean.
\end{enumerate}
\end{hbmethod}

The intersection in step~4 is valid because the two failure events have
total probability at most \(\delta/2+\delta/2\); this is the union
inequality again. The fixed-sample constructions of
\cref{sec:p15-bounded-concentration-learning} and the variance-sensitive
bounds of \cref{sec:p15-variance-dependent-concentration} can be
allocated over times in the same way, at the same logarithmic price.
Sequential tests with stopping and terminal decisions are developed in
\cref{sec:p15-sequential-evidence}, multiplicity over a fixed family in
\cref{sec:p15-multiple-testing}, and the stopping identities in
\cref{sec:p15-optional-stopping-and-random-walks}.

\emph{Sources.} The optimality statement, the Chebyshev comparison and the
normal-benchmark remarks are Hoeffding's
\cite[pp.~13--15 and~17]{Hoeffding1963BoundedSums}, and the comparison
of his three bounds is his Section~3
\cite[pp.~18--19]{Hoeffding1963BoundedSums}; the proofs supplied
for them are complete derivations from the results cited. The confidence
sequence definition, the stopping-time equivalence, the reuse failure,
the allocation and Wald comparisons, and the epoch construction are
derived here from \cref{thm:p15-conc-time-uniform},
\cref{thm:p15-conc-martingale}, \cref{thm:p15-doob-weak-maximal},
\cref{thm:p15-clt-oscillation} and the cited multiple-testing and
sequential-testing results; no further external presentation is used.
